\chapter*{\RL{ملخص}}
\addcontentsline{toc}{chapter}{Arabic Abstract}

\begin{RLtext}

\noindent يقدّم هذا التقرير الأعمال المنجزة في إطار مشروع نهاية الدراسات بشركة إس تي مايكروإلكترونيكس بتونس، والتي تتمحور حول قياس أداء المنظومة البرمجية STM32Cube مقارنةً بالمنظومات المنافسة (GigaDevice GD32، NXP، Texas Instruments، Renesas). يهدف المشروع إلى تزويد فريق التسويق التقني ببيانات مقارنة موضوعية وقابلة للتكرار حول البصمة الذاكرية (فلاش وRAM) وأداء التنفيذ ونضج الطبقات البرمجية الوسيطة.

صمّمنا منهجية قياس أداء مهيكلة تعتمد على سيناريوهات متطابقة تُنفَّذ على كلتا المنصتين المقارنتين بنفس مستوى التحسين وتردد الساعة والسلوك الملاحظ. طوّرنا بالتوازي أداتين متكاملتين: تطبيق \LR{MCU Workflow Studio} لتحليل البصمة الذاكرية عبر مقارنة ملفات الربط وتصنيف الرموز حسب الطبقة المعمارية، ونظام \LR{MCU Benchmark Copilot Agent} المدمج في بيئة \LR{VS Code} لتسريع إنشاء ونقل وتصحيح المعايير مع فرض التكافؤ الدلالي.

كشفت نتائج المقارنة بين ST وGD32 --- على مستوى مكدس USB (سيناريوهات HID وCDC) وتكامل FreeRTOS (سيناريوهات المعالجة الأساسية والتعاونية والاستباقية) --- عن مقايضة معمارية واضحة: يستهلك STM32 ذاكرة ROM أكثر (35--88\%) بسبب طبقة HAL المعيارية، لكنه يستخدم ذاكرة RAM أقل بكثير (57--65\% لـUSB) ويوفر أداء تهيئة أسرع بنسبة 61,7\%. قُدّمت مقترحات تحسين مستهدفة، أبرزها تحسين وحدة HAL RCC وتمكين الاختيار الدقيق لميزات USB عند التجميع.

\end{RLtext}

\noindent\rule[2pt]{\textwidth}{0.5pt}

\begin{RLtext}

{\textbf{الكلمات المفتاحية:}}

\LR{STM32Cube}، قياس الأداء، البصمة الذاكرية، متحكم دقيق، تحليل ملف الربط، \LR{FreeRTOS}، \LR{USB}، \LR{GD32}، النقل بين المصنّعين، وكيل ذكاء اصطناعي.
\\

\end{RLtext}

\noindent\rule[2pt]{\textwidth}{0.5pt}